\documentclass[10pt,a4paper]{amsart}
\usepackage[margin=19mm]{geometry}
\usepackage{amsmath,amssymb,mathtools}
\usepackage[scr=rsfs,cal=euler]{mathalfa}
\usepackage{xfrac}
\usepackage[unicode,hidelinks,bookmarks=false]{hyperref}
\newcommand{\ie}{i.e.\ }
\newcommand{\cf}{cf.\ }
\newcommand{\resp}{respectively\ }
\newcommand{\Subsection}{\subsection}
\providecommand{\nobreakdash}{\nobreak\hbox{-}}
\setlength{\emergencystretch}{2em}
\multlinegap0pt
\usepackage{color}
\usepackage{ulem}
\newtheorem{theorem}{Theorem}[section]
\newtheorem{corollary}[theorem]{Corollary}
\newtheorem{example}[theorem]{Example}
\newtheorem{assumption}[theorem]{Assumption}
\newtheorem{hypothesis}[theorem]{Hypothesis}
\newtheorem{proposition}[theorem]{Proposition}
\newtheorem{lemma}[theorem]{Lemma}
\newtheorem{definition}[theorem]{Definition}
\newtheorem{remark}[theorem]{Remark}
\def\bRzero{\mathbb R_0}
\def\cA{\mathcal{A}}
\def\cB{\mathcal{B}}
\def\cC{\mathcal{C}}
\def\cD{\mathcal{D}}
\def\cE{\mathcal{E}}
\def\cF{\mathcal{F}}
\def\cG{\mathcal{G}}
\def\cH{\mathcal{H}}
\def\cI{\mathcal{I}}
\def\cJ{\mathcal{J}}
\def\cL{\mathcal{L}}
\def\cM{\mathcal{M}}
\def\cN{\mathcal{N}}
\def\cO{\mathcal{O}}
\def\cP{\mathcal{P}}
\def\cR{\mathcal{R}}
\def\cQ{\mathcal{Q}}
\def\cS{\mathcal{S}}
\def\cT{\mathcal{T}}
\def\cU{\mathcal{U}}
\def\cV{\mathcal{V}}
\def\cX{\mathcal{X}}
\def\cZ{\mathcal{Z}}
\def\bC{\mathbb{C}}
\def\bD{\mathbb{D}}
\def\bE{\mathbb{E}}
\def\bQ{\mathbb{Q}}
\def\bN{\mathbb{N}}
\def\bP{\mathbb{P}}
\def\bR{\mathbb{R}}
\def\bS{\mathbb{S}}
\def\bZ{\mathbb{Z}}
\newcommand{\fH}{\mathfrak{H}}
\newcommand{\les}{\lesssim}
\def\wP{\widetilde{\cP}}
\def\Pb{\widetilde{\cP}_b}
\def\e{\varepsilon}
\def\k{\kappa}
\def\s{\sigma}
\def\SD{\mathbb{S}_{\mathbb{D}}}
\def\N{\overline{N}}
\def\barD{\overline{\mathbb{D}}_0}
\def\oZ{\overline{Z}}
\def\blue{\color{blue}}
\def\bs{{\bf s}}
\def\bt{{\bf t}}
\def\bu{{\bf u}}
\def\bv{{\bf v}}
\def\bx{{\bf x}}
\def\by{{\bf y}}
\def\bz{{\bf z}}
\newcommand{\fm}{\mathfrak{m}}
\begin{document}
\title[Central Limit Theorem for Stochastic Nonlinear Heat Equation]{Central Limit Theorem for Stochastic Nonlinear Heat Equation with Pure-Jump L\'evy White Noise}
\author{Matis Le Gall, Jinxin Wang}
\date{}
\maketitle
\noindent\emph{Source and licence.} Central Limit Theorem for Stochastic Nonlinear Heat Equation with Pure-Jump L\'evy White Noise. Version arXiv:2609.29023v1 (CC BY 4.0). This material is distributed under the Creative Commons Attribution 4.0 International licence, http://creativecommons.org/licenses/by/4.0/, as recorded in the arXiv OAI licence metadata for this exact version and shown on the abstract page https://arxiv.org/abs/2609.29023v1. Full rights notice: licenses/arxiv-2609.29023-CC-BY-4.0.txt.\par\smallskip
\noindent\textit{Editorial scope.} Counted unit: the original \texttt{theorem} environment \texttt{key-th-D2u} at source lines 622--654 together with its complete proof at lines 656--906; the delivered document keeps source lines 150--1022 so that every referenced label and every citation resolves inside it. Native editable source is the paper's own concatenated TeX; the AMS wrapper is editorial formatting only. No publisher class, upstream script or unrelated artwork is redistributed.
\begin{equation}
\label{nonlinear}
\begin{cases}
\displaystyle
\frac{\partial u}{\partial t}(t,x)
=
\frac{1}{2}\frac{\partial^2 u}{\partial x^2}(t,x)
+\sigma(u(t,x))\dot L(t,x),
\quad t>0,\ x\in\bR,\\[0.3em]
u(0,x)=1,\quad x\in\bR,
\end{cases}
\end{equation}
where $\dot L$ denotes the {\em L\'evy white noise} on
$\bR_+\times\bR$, and $\sigma:\bR\to\bR$ is a deterministic globally
Lipschitz function.


\medskip


 Let $N$ be a Poisson random measure on
$\bR_+\times\bR\times\bR_0$, with intensity $dtdx\nu(dz)$, and let
$\widehat N$ be its compensated version, where $\bR_0=\bR\setminus\{0\}$.
The L\'evy white noise $L$ is defined by
\begin{equation}
\label{def-L}
L(\varphi)
=
\int_{\bR_+\times\bR\times\bR_0}
\varphi(t,x)z\,
\widehat N(dt,dx,dz),
\qquad
\varphi\in L^2(\bR_+\times\bR).
\end{equation}
 For $p\ge1$, set $m_p:=\int_{\bR_0}|z|^p\nu(dz)$. Throughout, we
assume that $0<m_2<\infty$.


\medskip


\medskip

A predictable process $\{u(t,x):t\ge0,x\in\bR\}$ is called a mild solution of
\eqref{nonlinear} if
\begin{equation}
\label{solution}
u(t,x)
=
1+
\int_0^t\int_{\bR}
G_{t-s}(x-y)\sigma(u(s,y))L(ds,dy),
\end{equation}
where
\[
G_t(x)
=
\frac{1}{\sqrt{2\pi t}}
\exp\left(-\frac{x^2}{2t}\right),
\qquad t>0,\ x\in\bR,
\]
is the heat kernel. Existence for nonlinear stochastic heat equations driven by general
L\'evy noise was established in \cite{C17}, while path properties and
intermittency were studied in \cite{CDH19} and \cite{CK19}, respectively.
 These studies allow for more general L\'evy noises, for which the
finite-variance condition $m_2<\infty$ need not hold. For the finite-variance L\'evy white noise considered here, existence
and uniqueness of the random-field solution follow from
\cite{BN17}.
Throughout the paper we assume that $\sigma(1)\neq0$.
Otherwise, $u(t,x)\equiv1$ is the unique solution of
\eqref{nonlinear}, and the fluctuation problem is trivial. Since the stochastic integral in \eqref{solution} is centered,
$\bE[u(t,x)]=1$.
In the present paper,  our object of study is the centered spatial average
\begin{equation}
\label{FR-new}
F_R(t)
:=
\int_{-R}^{R}(u(t,x)-1)\,dx,
\qquad R>0.
\end{equation}


\medskip


When $\sigma(u)=u$, equation \eqref{nonlinear} becomes the
{\em parabolic Anderson model} (pAm) driven by L\'evy white noise. Spatial ergodicity, the limiting
covariance, a QCLT, and an FCLT for \eqref{FR-new} were established for that
model in \cite{BLW26}. The present work shows that this Gaussian fluctuation
picture persists for a globally Lipschitz nonlinear coefficient.



\medskip



A standard way to prove Gaussian fluctuations for spatial averages of SPDEs is
the Malliavin--Stein approach.  The key input  is a
$p$-moment estimate for the Malliavin derivatives of the solution.
Compared with the pAm analysis in \cite{BLW26}, the nonlinear coefficient
$\sigma$ does not change the structure of the first Malliavin derivative
estimate. More precisely, we prove that for any $p\in[2,3)$ such that
$m_p<\infty$,
\begin{equation}
\label{Du}
\big\|D_{r,y,z}u(t,x)\big\|_p
\le
C_{T,p}|z|\,g_{t-r}^{(p)}(x-y).
\end{equation}
 where $g_s^{(p)}(x)
:=
G_s(x)+G_s(x)^{2/p}$. The restriction $p<3$ comes from the integrability of the heat
kernel. For the stochastic nonlinear wave equation driven by
L\'evy white noise, Balan and Zheng \cite{BZ26} use a Poisson
Malliavin--Stein bound involving only first Malliavin derivatives.
Their proof relies on estimates for the fourth moment of the
first Malliavin derivative, so their argument cannot be applied
directly in our setting.




\medskip



We therefore use a second-order Poincar\'e inequality on the Poisson
space, which requires an estimate for $D^2u(t,x)$. This is where the nonlinear
coefficient creates a new minimum-type term. For $r_1<r_2$, the pAm
contribution has the product form
\[
|z_1z_2|\,
g_{t-r_2}^{(p)}(x-y_2)
g_{r_2-r_1}^{(p)}(y_2-y_1),
\]
whereas the nonlinear coefficient produces the additional term
\begin{equation}
\label{min-term}
\min\left\{
|z_1|g_{t-r_1}^{(p)}(x-y_1),
|z_2|g_{t-r_2}^{(p)}(x-y_2)
\right\}.
\end{equation}
The product term is already present in \cite{BLW26}. The minimum term
has no analogue in the pAm, and its control is the main technical contribution
of this paper.


\medskip

 For $T>0$, let $D[0,T]$ denote the space of c\`adl\`ag functions on
$[0,T]$. We write $J_1$ for the Skorokhod topology and $U$ for the topology of
uniform convergence. Our main results are as follows.


\begin{theorem}[Quantitative CLT]
\label{QCLT}
Assume that there exists $p\in(1,\frac32)$ such that
\begin{equation}
\label{mp-m2p}
m_1+m_{2p}<\infty.
\end{equation}
Then, for every $t>0$,
\[
{\rm dist}\left(\frac{F_R(t)}{\sigma_R(t)},Z\right)
\le C_t R^{-\left(1-\frac1p\right)} \qquad \text{for } R\ge1.
\]
Here $C_t>0$ depends on $t$, $Z$ is a standard normal random variable,
$\sigma_R^2(t)=\operatorname{Var}(F_R(t))$, and ${\rm dist}$ denotes any one
of the distances $d_W$, $d_{FM}$, and $d_K$, namely the $1$-Wasserstein,
Fortet--Mourier, and Kolmogorov distances, respectively.
\end{theorem}

\begin{theorem}[Functional CLT]
\label{FCLT}
Under the hypotheses of Theorem~\ref{QCLT}, for every $R>0$, the process
$\{F_R(t)\}_{t\geq0}$ has a c\`adl\`ag modification (denoted also $F_R$).
Moreover, for any $T>0$,
\[
\frac{1}{\sqrt R}F_R(\cdot)\xrightarrow{d}\cG(\cdot)
\quad \mbox{in $(D[0,T],J_1)$ as $R\to\infty$},
\]
where $\{\cG(t)\}_{t\geq0}$ is a centered continuous Gaussian process with covariance
\[
\bE[\cG(t)\cG(s)]
=\Sigma_{t,s}
:=2m_2\int_0^{t\wedge s}\bE|\sigma(u(r,0))|^2\,dr.
\]
The convergence also holds in $(D[0,T],U)$.
\end{theorem}


Section~2 establishes the required moment bounds for $u$, $Du$, and
$D^2u$. Section~3 applies these bounds to the spatial averages and proves
Theorems~\ref{QCLT} and \ref{FCLT}.


\section{Key estimates on Malliavin derivatives}
\subsection{Poisson Malliavin derivative and chain estimates}
Let $({\bf Z},\cZ,\fm)$ be the measure space
\[
({\bf Z},\cZ,\fm)
=\big(\bR_+\times\bR\times\bR_0,
\cB(\bR_+)\otimes\cB(\bR)\otimes\cB(\bR_0),
{\rm Leb}\otimes{\rm Leb}\otimes\nu\big),
\]
 and set $\cH=L^2({\bf Z},\cZ,\fm)$. For $q\ge1$, we write
$\|X\|_q=(\bE|X|^q)^{1/q}$. Constants denoted by $C$ or $C_{T,p}$ may
change from line to line. For Malliavin calculus with respect to the
compensated Poisson random measure $\widehat N$, we refer to \cite{NN18}.

\medskip

We recall the add-one Malliavin derivative. Let ${\bf N}_{\bf Z}$ be the space of all $\sigma$-finite counting measures
on $({\bf Z},\cZ)$. If $F$ is a real-valued $\cF^N$-measurable random
variable, then there exists a measurable function
$f:{\bf N}_{\bf Z}\to\bR$ such that $F=f(N)$. For $\xi\in{\bf Z}$, the
add-one operator is defined by
\begin{equation}
\label{add-one}
D_\xi^+F
:=
f(N+\delta_\xi)-f(N),
\end{equation}
whenever the right-hand side is well defined. 

It is known that the add-one operator $D^+$ coincides with the Malliavin
derivative operator $D$ on ${\rm dom}(D)$. Hence, we simply
write $D_\xi F$ instead of $D_\xi^+F$. For $\xi_1,\xi_2\in{\bf Z}$, the
second-order difference is given by
\begin{equation}
\label{second-add-one}
D^2_{\xi_1,\xi_2}F
=
f(N+\delta_{\xi_1}+\delta_{\xi_2})
-
f(N+\delta_{\xi_1})
-
f(N+\delta_{\xi_2})
+
f(N),
\end{equation}
whenever the right-hand side is well defined.

The following elementary chain estimates will be used repeatedly.

\begin{lemma}
\label{lem:chain-estimates}
Let $\phi:\bR\to\bR$ be a Lipschitz function with Lipschitz constant
${\rm Lip}(\phi)$. Then, whenever the following differences are well defined,
\begin{equation}
\label{D-chain}
|D_\xi\phi(F)|
\le
{\rm Lip}(\phi)|D_\xi F|,
\end{equation}
and
\begin{equation}
\label{D2-chain}
|D^2_{\xi_1,\xi_2}\phi(F)|
\le
{\rm Lip}(\phi)|D^2_{\xi_1,\xi_2}F|
+
2{\rm Lip}(\phi)
\min\big\{
|D_{\xi_1}F|,
|D_{\xi_2}F|
\big\}.
\end{equation}
\end{lemma}

\begin{proof}
The first estimate follows directly from the Lipschitz property and
\eqref{add-one}. For the second estimate, we have
\begin{align*}
D^2_{\xi_1,\xi_2}\phi(F)
&=
\phi\big(f(N+\delta_{\xi_1}+\delta_{\xi_2})\big)
-
\phi\big(f(N+\delta_{\xi_1})\big)
-
\phi\big(f(N+\delta_{\xi_2})\big)
+
\phi(f(N)) \\
&=
\phi\left(
F+D_{\xi_1}F+D_{\xi_2}F+D^2_{\xi_1,\xi_2}F
\right)
-
\phi\left(F+D_{\xi_1}F\right)
-
\phi\left(F+D_{\xi_2}F\right)
+
\phi(F).
\end{align*}
We use the elementary Lipschitz estimate
\[
|\phi(x+a+b+c)-\phi(x+a)-\phi(x+b)+\phi(x)|
\le
{\rm Lip}(\phi)
\big(
|c|+2\min\{|a|,|b|\}
\big).
\]
Applying this estimate with $x=F$,  $a=D_{\xi_1}F$, $b=D_{\xi_2}F$ and $c=D^2_{\xi_1,\xi_2}F$ gives the desired bound.

\end{proof}













\subsection{Malliavin derivative estimates}
 We first show that the solution has finite $p$-th moments for
$p\in[2,3)$ with $m_p<\infty$, and then estimate its first and second
Malliavin derivatives. The arguments extend the corresponding pAm estimates,
so we emphasize the modifications caused by the nonlinear coefficient. We write $L_\sigma$ for a Lipschitz constant of $\sigma$. 
Throughout this section, for $p\ge2$, set
\[
H_t^{(p)}(x)
=
G_t^2(x)+G_t^p(x) \qquad \text{and}\qquad g_t^{(p)}(x)
=
G_t(x)+G_t^{2/p}(x).
\]



We begin with the moment estimate.
\begin{theorem}
\label{thm-u-pth-moment-nonlinear}
Let $p\in [2,3)$ be such that $m_p<\infty$. Let $u$ be the mild solution of
\eqref{nonlinear}. Then, for any $T>0$,
\begin{equation}
\label{u-pth-moment}
    \sup_{(t,x)\in [0,T]\times \bR} \bE|u(t,x)|^p <\infty .
\end{equation}
\end{theorem}

\begin{proof}
Let $(u_n)_{n\geq 0}$ be the sequence of Picard iterations, defined by
$u_0(t,x)=1$ and
\begin{equation}
\label{Picard}
u_{n+1}(t,x)=1+\int_0^t \int_{\bR} G_{t-s}(x-y) \sigma (u_n(s,y)) L(ds,dy),
\quad n\geq 0.
\end{equation}
By \cite[Lemma~2.5]{BLW26}, the linear growth of $\sigma$, and
$\int_0^T\int_{\bR}H_s^{(p)}(y)\,dy\,ds<\infty$,
\[
\bE|u_{n+1}(t,x)|^p
\le
C_{T,p}
+C_{T,p}
\int_0^t\int_{\bR}
H_{t-s}^{(p)}(x-y)\bE|u_n(s,y)|^p\,dy\,ds .
\]
Iterating this renewal inequality as in \cite[Theorem~2.7]{BLW26}
yields a bound uniform in $n$. Passing to the limit in the Picard sequence
proves \eqref{u-pth-moment}.
\end{proof}


\begin{theorem}
\label{key-th-Du}
\begin{itemize}
\item[(i)] For every $t\geq0$ and $x\in\bR$, $u(t,x)\in{\rm dom}(D)$.
\item[(ii)] Let $p\in[2,3)$ satisfy $m_p<\infty$. For any
$0\le r<t\le T$, $x,y\in\bR$, and $z\in\bR_0$,
\begin{equation}
\label{Du-pth-moment}
 \| D_{r,y,z}u(t,x)\|_p\le C_{T,p}|z|g_{t-r}^{(p)}(x-y),
\end{equation}
where $C_{T,p}>0$ is non-decreasing in $T$.
\end{itemize}
\end{theorem}

\begin{proof}
By induction, $u_n(t,x)\in{\rm dom}(D)$ for every $n\ge0$.
Let $\xi=(r,y,z)$. From the recurrence relation \eqref{Picard} and the Heisenberg commutation principle, it follows that
\begin{equation*}
D_\xi u_{n+1}(t,x)=  zG_{t-r}(x-y)\sigma(u_n(r,y)) + \int_r^t\int_{\bR}
G_{t-s_1}(x-y_1)D_\xi\sigma(u_n(s_1,y_1))L(ds_1,dy_1).
\end{equation*}
 By \cite[Lemma~2.5]{BLW26}, \eqref{u-pth-moment}, and
\eqref{D-chain},
\begin{align*}
 \bE\left|D_{\xi}u_{n+1}(t,x)\right|^p \le &  2^{p-1}C_{T,p}|z|^pG_{t-r}^p(x-y)\\
 &+2^{p-1}A_{T,p}L_\sigma^p
\int_r^t\int_{\bR}
H_{t-s_1}^{(p)}(x-y_1)
\bE\left|D_{\xi}u_n(s_1,y_1)\right|^p
dy_1ds_1.
\end{align*}
The renewal argument in \cite[Theorem~3.1]{BLW26} now gives a bound
uniform in $n$. The closability of $D$ and convergence of the Picard sequence
yield (i) and \eqref{Du-pth-moment}.
\end{proof}



\medskip




For $k\ge1$ and $0\le r<t$, define
$T_k(r,t)=\{(t_1,\ldots,t_k):r<t_1<\cdots<t_k<t\}$. We use the
conventions $t_{k+1}=t$, $x_{k+1}=x$,
$\pmb{t}=(t_1,\ldots,t_k)$, and $\pmb{x}=(x_1,\ldots,x_k)$.
In \cite{BLW26},  it is proved that if $p\in[2,3)$ and
$T>0$, then for any $0\le r<t\le T$ and $x,y\in\bR$,
\begin{align}
\label{Gp-bound}
\sum_{k=1}^{\infty}
C_{T,p}^k
\int_{T_k(r,t)}
\int_{\bR^k}
\prod_{i=1}^{k}
H_{t_{i+1}-t_i}^{(p)}(x_{i+1}-x_i)
G_{t_1-r}^p(x_1-y)
\,d\pmb{x}\,d\pmb{t}
\le
C_{T,p}H_{t-r}^{(p)}(x-y).
\end{align}
The proof of \eqref{Gp-bound} splits into two cases:
the case in which all factors are of type $G^p$, and the
case in which at least one factor is of type $G^2$.
The following estimate is an extension of
\eqref{Gp-bound}. The difference is that the initial kernel
$G_{t_1-r}^p$ is replaced by $H_{t_1-r}^{(p)}$.


\begin{lemma}
Let $p\in[2,3)$ and $T>0$. Then, for any $0\le r<t\le T$ and
$x,y\in\bR$,
\begin{align}
\label{GpG2-bound}
\sum_{k=1}^{\infty}
C_{T,p}^k
\int_{T_k(r,t)}
\int_{\bR^k}
\prod_{i=1}^{k}
H_{t_{i+1}-t_i}^{(p)}(x_{i+1}-x_i)
H_{t_1-r}^{(p)}(x_1-y)
\,d\pmb{x}\,d\pmb{t}
\le
C_{T,p}
H_{t-r}^{(p)}(x-y).
\end{align}
\end{lemma}

\begin{proof}
 Since $H_{t_1-r}^{(p)}=G_{t_1-r}^p+G_{t_1-r}^2$, the first
contribution is bounded by \eqref{Gp-bound}. For the second, expand the product
of the $H^{(p)}$ kernels. Every resulting term contains a $G^2$ factor, so the
second case in the proof of \eqref{Gp-bound} applies.
\end{proof}






\begin{theorem}
\label{key-th-D2u}
If $m_1<\infty$, then
\begin{itemize}
\item[(i)] for every $t>0$ and $x\in\bR$, $u(t,x)\in{\rm dom}(D^2)$;
\item[(ii)] if $p\in[2,3)$ and $m_p<\infty$, then for any
$0\le t\le T$, $x\in\bR$, and
$\xi_i=(r_i,y_i,z_i)\in{\bf Z}$ with $r_i\in[0,t]$, $i=1,2$,
\begin{equation}
\label{D2u-nonlinear-bound}
\begin{aligned}
\|D^2_{\xi_1,\xi_2}u(t,x)\|_p
&\le
C_{T,p}|z_1z_2|
\begin{cases}
g_{t-r_2}^{(p)}(x-y_2)
g_{r_2-r_1}^{(p)}(y_2-y_1),
& \text{if } r_1<r_2, \\[0.3em]
g_{t-r_1}^{(p)}(x-y_1)
g_{r_1-r_2}^{(p)}(y_1-y_2),
& \text{if } r_2<r_1,
\end{cases}
\\
&\quad+
C_{T,p}
\min\left\{
|z_1|g_{t-r_1}^{(p)}(x-y_1),
|z_2|g_{t-r_2}^{(p)}(x-y_2)
\right\}.
\end{aligned}
\end{equation}
\end{itemize}
\end{theorem}

\begin{proof}
We prove the estimate for $0\le r_1<r_2<t\le T$. The case
$r_2<r_1$ follows by symmetry, and the diagonal $r_1=r_2$ is
$\fm^{\otimes2}$-null. By induction, $u_n(t,x)\in{\rm dom}(D^2)$ for every
$n\ge0$.
For the Picard approximations, applying $D_{\xi_1}$ to the equation for
$D_{\xi_2}u_{n+1}(t,x)$ gives
\begin{align}
\label{D2u-picard-nonlinear}
D^2_{\xi_1,\xi_2}u_{n+1}(t,x)
=&
z_2G_{t-r_2}(x-y_2)
D_{\xi_1}\sigma(u_n(r_2,y_2))
\nonumber\\
&+
\int_{r_2}^t\int_{\bR}
G_{t-s}(x-y)
D^2_{\xi_1,\xi_2}\sigma(u_n(s,y))
L(ds,dy).
\end{align}
Applying Rosenthal's inequality with
$\Phi(s,y)=D^2_{\xi_1,\xi_2}\sigma(u_n(s,y))$, we obtain
\begin{equation}
    \label{D2u-0}
\begin{aligned}
 \bE\left|D^2_{\xi_1,\xi_2}u_{n+1}(t,x)\right|^p
&\le
2^{p-1}
\bigg\{
|z_2|^pG_{t-r_2}^p(x-y_2)
\bE\left|D_{\xi_1}\sigma\big(u_n(r_2,y_2)\big)\right|^p\\
&\quad+
A_{T,p}
\int_{r_2}^t\int_{\bR}
H_{t-s}^{(p)}(x-y)
\bE\left|
D^2_{\xi_1,\xi_2}\sigma(u_n(s,y))
\right|^p
dy ds
\bigg\}.
\end{aligned}
\end{equation}
By \eqref{D-chain} and 
 \eqref{Du-pth-moment},
\begin{equation}
    \label{D2u-1}
\bE\left|D_{\xi_1}\sigma\big(u_n(r_2,y_2)\big)\right|^p\le
L_\sigma^p C_{T,p}^p
|z_1|^p
H_{r_2-r_1}^{(p)}(y_2-y_1). 
\end{equation}
By \eqref{D2-chain},
\begin{equation}
\label{D2u-2}
\begin{aligned}
\bE\left|D^2_{\xi_1,\xi_2}\sigma\big(u_n(s,y)\big)\right|^p
&\le
C_pL_\sigma^p
\bE\left|D^2_{\xi_1,\xi_2}u_n(s,y)\right|^p
+
C_pL_\sigma^p
\bE\left[
\min_{i=1,2}
\left|D_{\xi_i}u_n(s,y)\right|^p
\right]  \\
&\le
C_pL_\sigma^p
\bE\left|D^2_{\xi_1,\xi_2}u_n(s,y)\right|^p
+
C_{T,p}L_\sigma^p
\min_{i=1,2}
\left\{
|z_i|^pH_{s-r_i}^{(p)}(y-y_i)
\right\}.
\end{aligned}
\end{equation}
Combining \eqref{D2u-0} with \eqref{D2u-1} and \eqref{D2u-2}, we obtain
\begin{equation}
\label{D2-renewal-nonlinear}
\begin{aligned}
\bE\left|D^2_{\xi_1,\xi_2}u_{n+1}(t,x)\right|^p
&\le
C_{T,p}|z_1z_2|^p
G_{t-r_2}^p(x-y_2)
H_{r_2-r_1}^{(p)}(y_2-y_1)
\\
&\quad+
C_{T,p}
\int_{r_2}^t\int_{\bR}
H_{t-s}^{(p)}(x-y)
\bE\left|D^2_{\xi_1,\xi_2}u_n(s,y)\right|^p
dy ds
\\
&\quad+
C_{T,p}
\int_{r_2}^t\int_{\bR}
H_{t-s}^{(p)}(x-y)
\min_{i=1,2}
\left\{
|z_i|^pH_{s-r_i}^{(p)}(y-y_i)
\right\}
dy ds .
\end{aligned}
\end{equation}
 For $n\ge2$, iterating \eqref{D2-renewal-nonlinear} and using
$D^2_{\xi_1,\xi_2}u_1=0$, we obtain
\begin{align*}
\bE\left|D^2_{\xi_1,\xi_2}u_n(t,x)\right|^p&\le C_{T,p}|z_1z_2|^p
H_{r_2-r_1}^{(p)}(y_2-y_1) \left[ G_{t-r_2}^p(x-y_2)+ \sum_{k=1}^{n-2}
I_1^{(k)}\right]+\sum_{k=1}^{n-1}
 I_2^{(k)}
\end{align*}
where
\begin{equation*}
I_1^{(k)}= C_{T,p}^k
\int_{T_k(r_2,t)}
\int_{\bR^k}
\prod_{i=1}^{k}
H_{t_{i+1}-t_i}^{(p)}(x_{i+1}-x_i)\cdot
G_{t_1-r_2}^p(x_1-y_2)
d\pmb{x}d\pmb{t}  
\end{equation*}
and 
\begin{equation*}
I_2^{(k)}= C_{T,p}^k
\int_{T_k(r_2,t)}
\int_{\bR^k}
\prod_{i=1}^{k}
H_{t_{i+1}-t_i}^{(p)}(x_{i+1}-x_i) \cdot
\min_{i=1,2}\left\{
|z_i|^pH_{t_1-r_i}^{(p)}(x_1-y_i)
\right\}
d\pmb{x}d\pmb{t},
\end{equation*}
By \eqref{Gp-bound},
\begin{equation}
\label{D2-source-iteration-bound}
G_{t-r_2}^p(x-y_2)
+
\sum_{k=1}^{n-2} I_1^{(k)}
\le
C_{T,p}H_{t-r_2}^{(p)}(x-y_2).
\end{equation}


For each fixed $j\in\{1,2\}$, since $\min_{i=1,2}
\left\{
|z_i|^pH_{t_1-r_i}^{(p)}(x_1-y_i)
\right\}
\le
|z_j|^pH_{t_1-r_j}^{(p)}(x_1-y_j)$, and since $T_k(r_2,t)\subset T_k(r_j,t)$ for $j=1,2$, we obtain
\begin{align*}
\sum_{k=1}^{n-1} I_2^{(k)}
&\le
\sum_{k=1}^{n-1}
C_{T,p}^k
\int_{T_k(r_j,t)}
\int_{\bR^k}
\prod_{i=1}^{k}
H_{t_{i+1}-t_i}^{(p)}(x_{i+1}-x_i)
|z_j|^pH_{t_1-r_j}^{(p)}(x_1-y_j)
d\pmb{x}d\pmb{t}.
\end{align*}
Therefore, by \eqref{GpG2-bound}
\begin{equation*}
\sum_{k=1}^{n-1} I_2^{(k)}
\le
C_{T,p}
|z_j|^pH_{t-r_j}^{(p)}(x-y_j),
\qquad j=1,2.
\end{equation*}
Taking the minimum over $j=1,2$, we get
\begin{equation}
\label{D2-min-iteration-bound}
\sum_{k=1}^{n-1} I_2^{(k)}
\le
C_{T,p}
\min\left\{
|z_1|^pH_{t-r_1}^{(p)}(x-y_1),
|z_2|^pH_{t-r_2}^{(p)}(x-y_2)
\right\}.
\end{equation}





Using \eqref{D2-source-iteration-bound} and
\eqref{D2-min-iteration-bound}, we obtain, uniformly in $n\ge2$,
\begin{align}
\label{D2u-G-nonlinear}
\bE\left|D^2_{\xi_1,\xi_2}u_n(t,x)\right|^p
&\le
C_{T,p}|z_1z_2|^p
H_{t-r_2}^{(p)}(x-y_2)
H_{r_2-r_1}^{(p)}(y_2-y_1)
\nonumber\\
&\quad+
C_{T,p}
\min\left\{
|z_1|^pH_{t-r_1}^{(p)}(x-y_1),
|z_2|^pH_{t-r_2}^{(p)}(x-y_2)
\right\}.
\end{align}



Taking $p=2$ in \eqref{D2u-G-nonlinear} and integrating with respect to
$\fm(d\xi_1)\fm(d\xi_2)$ gives
\[
\sup_{n\ge2}
\bE\|D^2u_n(t,x)\|_{\cH^{\otimes2}}^2<\infty.
\]
The product term is integrable because $m_2<\infty$. For the minimum
term, use

\[
\min\{|z_1|^2A,|z_2|^2B\}\le |z_1z_2|A^{1/2}B^{1/2},
\qquad (H_s^{(2)})^{1/2}=\sqrt2\,G_s.
\]
Its contribution is bounded by
\[
\begin{aligned}
&\int_0^t\int_0^t\int_{\bR^2}\int_{\bR_0^2}
\min\left\{|z_1|^2H_{t-r_1}^{(2)}(x-y_1),
|z_2|^2H_{t-r_2}^{(2)}(x-y_2)\right\}
\nu(dz_1)\nu(dz_2)dy_1dy_2dr_1dr_2 \\
&\qquad\le
2m_1^2
\int_0^t\int_0^t
\left(\int_{\bR}G_{t-r_1}(x-y_1)\,dy_1\right)
\left(\int_{\bR}G_{t-r_2}(x-y_2)\,dy_2\right)
dr_1dr_2  \\
&\qquad=
2m_1^2t^2
<\infty .
\end{aligned}
\]
 Thus the minimum term is integrable because $m_1<\infty$.





\medskip
Apply \cite[Lemma~2.1]{BLW26} with $F_n=u_n(t,x)$ and $F=u(t,x)$
for fixed $(t,x)\in\bR_+\times\bR$. It follows that
$u(t,x)\in{\rm dom}(D^2)$ and that $D^2u_n(t,x)$ converges weakly to
$D^2u(t,x)$ in $L^2(\Omega;\cH^{\otimes2})$. This proves (i). Part (ii)
follows from \cite[Lemma~A.1]{BS26}.
\end{proof}



\section{Gaussian fluctuations of the spatial averages}
\label{sec:Gaussian-fluctuations}

We now prove the Gaussian fluctuation results. Ergodicity, the limiting
covariance, and the FCLT follow the pAm arguments after inserting the nonlinear
moment bounds. The QCLT additionally requires control of the minimum term in
the second Malliavin derivative estimate.

For $0\le r<t$, $R>0$ and $y\in\bR$, set
\begin{equation}
\label{def-phi}
\varphi_{t,R}(r,y)
=
\int_{-R}^{R}G_{t-r}(x-y)\,dx,
\qquad
\varphi_{t,R}(y)=\varphi_{t,R}(0,y).
\end{equation}

\subsection{Ergodicity and limiting covariance}
\begin{theorem}
\label{ergodic-cov}
\begin{itemize}
\item[(i)] For every $t>0$, the random field $\{u(t,x)\}_{x\in\bR}$ is
strictly stationary and ergodic. Consequently,
\[
\lim_{R\to\infty}\frac{1}{R}F_R(t)=0
\qquad
\text{a.s. and in }L^2(\Omega).
\]

\item[(ii)] For every $s,t>0$,
\begin{equation}
\label{limit-cov}
\lim_{R\to\infty}
\frac{1}{R}\bE[F_R(t)F_R(s)]
=
\Sigma_{t,s}
:=
2m_2\int_0^{t\wedge s}
\bE\big|\sigma(u(r,0))\big|^2\,dr.
\end{equation}
In particular,
\[
\sigma_R^2(t):=\operatorname{Var}(F_R(t))
\sim
\Sigma_{t,t}R,
\qquad R\to\infty.
\]
\end{itemize}
\end{theorem}

\begin{proof}
 Part~(i) follows as in \cite[Theorem~1.1]{BZ24}: stationarity follows
from the spatial homogeneity of $L$, and ergodicity from
\cite[Lemma~4.2]{BZ24} and Theorem~\ref{key-th-Du} with $p=2$.

As in \cite[proof of Theorem~1.1(ii)]{BLW26}, stochastic Fubini's
theorem, the It\^o isometry, and stationarity give
\[
\bE[F_R(t)F_R(s)]
=
m_2\int_0^{t\wedge s}
\bE|\sigma(u(r,0))|^2
\int_{\bR}
\varphi_{t,R}(r,y)\varphi_{s,R}(r,y)\,dy\,dr,
\]
 and, for $r<t\wedge s$,
\[
\frac1R\int_{\bR}
\varphi_{t,R}(r,y)\varphi_{s,R}(r,y)\,dy
\longrightarrow 2,
\qquad R\to\infty.
\]
Dominated convergence yields \eqref{limit-cov}.
\end{proof}

\subsection{Quantitative central limit theorem}
We first introduce some notation. For \(t>0\), \(R>0\),
\(p>0\) and \(y\in\bR\), define
\[
K_t^{(p)}(y)
=
G_t(y)+G_{pt}(y),
\qquad
\Phi_{t,R}^{(p)}(y)
=
\int_{-R}^{R}K_t^{(p)}(x-y)\,dx .
\]

We use the estimates from \cite[Lemmas~4.3--4.5]{BLW26}, collected
below for convenience.
\begin{lemma}
Fix $T>0$. For every $0<t\le T$, $R>0$ and $p>0$,
\begin{align}
\label{Phi-bound}
\Phi_{t,R}^{(p)}(y)
&\le 2,
 \\
\label{K-power-bound}
\big(K_t^{(p)}(x)\big)^\theta
&\lesssim K_{t/\theta}^{(p)}(x),
\qquad 0<\theta<1,\\
\label{K-power-bound-large}
\big(K_t^{(p)}(x)\big)^\theta
&\lesssim t^{-\frac{\theta-1}{2}}K_{t/\theta}^{(p)}(x),
\qquad \theta>1.
\end{align}
\end{lemma}


The following result is the nonlinear analogue of
\cite[Lemma~4.2]{BLW26}. It follows by integrating the estimates of
Theorems~\ref{key-th-Du} and \ref{key-th-D2u} over $[-R,R]$, so we omit its proof.


\begin{thebibliography}{99}
\bibitem[BLW26]{BLW26} Balan, R.M., Le Gall, M. and Wang, J. (2026). Gaussian fluctuations for the parabolic Anderson model with L\'evy white noise. Preprint, arXiv:2607.02742.
\bibitem[BN17]{BN17} Balan, R.M. and Ndongo, C.B. (2017). Malliavin differentiability of solutions of SPDEs with L´evy white noise. {\em International Journal of Stochastic Analysis.} %{\bf 87}, 363-381.
\bibitem[BS26]{BS26} Balan, R.M. and Stephenson, W.D. (2026). Gaussian fluctuations for hyperbolic Anderson model with L´evy colored noise. Preprint arXiv:2602.23137.
\bibitem[BZ24]{BZ24} Balan, R.M. and Zheng, G. (2024). Hyperbolic Anderson model with L\'evy white noise: spatial ergodicty and fluctuation. {\em Trans. AMS} {\bf 377}, 4171-4221.
\bibitem[BZ26]{BZ26} Balan, R.M. and Zheng, G. (2026). Central limit theorem for stochastoc nonlinear wave equation with pure-jump L\'evy white noise. {\em Electr. J. Probab.} {\bf 31}, no. 70, 1-46.
\bibitem[C17]{C17} Chong, C. (2017). Stochastic PDEs with heavy-tailed noise. {\em Stochastic Process. Appl.} {\bf 127}, 2262--2280.
\bibitem[CDH19]{CDH19} Chong, C., Dalang, R.C. and Humeau, T. (2019). Path properties of the solution to the stochastic heat equation with L\'evy noise. {\em Stoch. Partial Differ. Equ. Anal. Comput.} {\bf 7}, 123--168.
\bibitem[CK19]{CK19} Chong, C. and Kevei, P. (2019). Intermittency for the stochastic heat equation with L\'evy noise. {\em Ann. Probab.} {\bf 47}, 1911--1948.
\bibitem[NN18]{NN18} Nualart, D. and Nualart, E. (2018). {\em Introduction to Malliavin Calculus}. Cambridge University Press, Cambridge.
\end{thebibliography}
\end{document}
